%%%%%%%%%%%%%%%%%%%%%%%%%%%%%%%%%%%%%%%%%%%%%%%%%%%%%%%%%%%%%%%%%%%%%%%%%%%%%%%%%%%%%%%%%%
% RÉSUMÉ
%%%%%%%%%%%%%%%%%%%%%%%%%%%%%%%%%%%%%%%%%%%%%%%%%%%%%%%%%%%%%%%%%%%%%%%%%%%%%%%%%%%%%%%%%%

\chapter*{Résumé}
\addstarredchapter{Résumé}

\initial{L}’industrialisation des modèles d’intelligence artificielle représente aujourd’hui un enjeu majeur pour les entreprises souhaitant déployer des solutions performantes en environnement de production. Si les modèles développés en phase expérimentale offrent généralement de bons résultats, leur passage à l’échelle est souvent limité par le volume croissant des données, les besoins en ressources de calcul et les contraintes liées à l’automatisation des traitements. Ce mémoire, réalisé au sein de \textbf{DESEMITS LLC}, porte sur la conception d’une architecture Cloud orientée MLOps destinée à accompagner le passage à l’échelle d’un modèle de vision par ordinateur utilisé pour la détection d’anomalies sur des cultures agricoles. L’objectif est de concevoir une infrastructure distribuée capable d’assurer un traitement efficace de grands volumes d’images tout en garantissant la scalabilité, la disponibilité et l’optimisation des coûts. La solution proposée repose sur les services \textbf{Amazon Web Services (AWS)}, notamment Amazon S3 pour le stockage distribué des données, Amazon EMR pour le calcul distribué et Apache Spark (PySpark) pour le traitement parallèle des images. Cette architecture permet de séparer les ressources de stockage et de calcul, d’améliorer les performances du pipeline de traitement et de faciliter son intégration dans une démarche MLOps. Les expérimentations réalisées montrent que l’architecture proposée répond efficacement aux exigences de montée en charge en offrant une solution modulaire, évolutive et adaptée au traitement massif des données. Elle constitue ainsi une base robuste pour le déploiement et l’industrialisation de modèles de vision par ordinateur dans un environnement Cloud.

\vspace{0.5cm}

Mots-clés : \textbf{Intelligence artificielle, MLOps, Cloud Computing, Amazon Web Services, Amazon EMR, Amazon S3, Apache Spark, Calcul distribué, Vision par ordinateur, Scalabilité.}

\clearpage